\chapter{Strong Form, Weak Form and Variational Principles}
\label{ch:strong-weak}

The finite element method does not attack the governing partial
differential equation directly. It works from an equivalent
\emph{integral} (weak, or variational) statement. This chapter derives
that statement and shows its three faces --- virtual work, weighted
residuals, and energy minimisation --- which coincide for
elastostatics. The presentation follows Hughes' canonical
(S)/(W)/(G) development~\citep{hughes2000}, with Zienkiewicz \&
Taylor~\citep{zienkiewicz2013} and Bathe~\citep{bathe2014}.

\section{The strong form (S)}

Let $\domain\subset\mathbb R^d$ have boundary
$\boundary=\GammaD\cup\GammaN$ with $\GammaD\cap\GammaN=\emptyset$. The
\textbf{strong form} of linear elastostatics is: find
$\disp:\domain\to\mathbb R^d$ such that
\begin{subequations}\label{eq:strong}
\begin{align}
  \divergence\stress + \bodyforce &= \vec 0 && \text{in }\domain
    &&\text{(equilibrium)},\\
  \stress &= \Dmat:\strain(\disp) && \text{in }\domain
    &&\text{(constitution)},\\
  \strain &= \sym\grad\disp && \text{in }\domain
    &&\text{(kinematics)},\\
  \disp &= \bar{\disp} && \text{on }\GammaD
    &&\text{(Dirichlet / essential)},\\
  \stress\,\vec n &= \bar{\traction} && \text{on }\GammaN
    &&\text{(Neumann / natural)}.
\end{align}
\end{subequations}
``Strong'' because it demands the PDE hold pointwise, which requires
$\disp$ to be twice differentiable. Real solutions --- near corners,
point loads, material interfaces --- often are not that smooth, and the
strong form is awkward to discretise. The weak form removes one order of
differentiability.

\section{Deriving the weak form (W)}

Introduce a \textbf{test function} (weighting function, or virtual
displacement) $\vec w$ that is arbitrary except that it vanishes where
the solution is already prescribed: $\vec w=\vec 0$ on $\GammaD$.
Multiply the equilibrium equation by $\vec w$ and integrate over the
body:
\begin{equation}
  \int_\domain \vec w\cdot(\divergence\stress+\bodyforce)\dd V = 0 .
\end{equation}
Integrate the stress-divergence term by parts (the divergence theorem):
\begin{equation}
  \int_\domain \vec w\cdot(\divergence\stress)\dd V
   = \int_{\boundary}\vec w\cdot(\stress\vec n)\dd A
   - \int_\domain \grad\vec w : \stress \dd V .
\end{equation}
Because $\stress$ is symmetric, $\grad\vec w:\stress
=\sym\grad\vec w:\stress=\strain(\vec w):\stress$. On the boundary,
$\vec w=\vec 0$ on $\GammaD$ kills that portion, and
$\stress\vec n=\bar{\traction}$ on $\GammaN$. Rearranging gives the
\textbf{weak form}:

\begin{tcolorbox}[colback=gray!5,colframe=gray!55!black,arc=2pt]
Find $\disp\in\mathcal S$ such that, for all $\vec w\in\mathcal V$,
\begin{equation}
  \underbrace{\int_\domain \strain(\vec w):\Dmat:\strain(\disp)\dd V}_{\text{internal virtual work}}
  \;=\;
  \underbrace{\int_\domain \vec w\cdot\bodyforce\dd V
   + \int_{\GammaN}\vec w\cdot\bar{\traction}\dd A}_{\text{external virtual work}} .
  \label{eq:weak}
\end{equation}
\end{tcolorbox}

\noindent In Voigt/matrix form the bilinear term reads
$\int_\domain (\grad_{\!s}\vec w)\transpose\,\Dmat\,(\grad_{\!s}\disp)\dd V$.
Equation~\eqref{eq:weak} is the \textbf{principle of virtual work}:
internal virtual work equals external virtual work for every admissible
virtual displacement.

\subsection{Function spaces}

The integrals in~\eqref{eq:weak} require only \emph{first} derivatives to
be square-integrable, i.e.\ finite strain energy. The natural setting is
the Sobolev space $H^1(\domain)$. Trial solutions and test functions
live in
\begin{align}
  \mathcal S &= \{\disp\in[H^1(\domain)]^d : \disp=\bar{\disp}\text{ on }\GammaD\}
    &&\text{(trial space)},\\
  \mathcal V &= \{\vec w\in[H^1(\domain)]^d : \vec w=\vec 0\text{ on }\GammaD\}
    &&\text{(test space)}.
\end{align}
The Dirichlet condition is built into $\mathcal S$ --- hence
``essential'' --- while the Neumann condition emerged from the boundary
integral --- hence ``natural.'' This weaker regularity requirement is
exactly what lets us use simple piecewise-polynomial (finite element)
approximations, which are continuous but only piecewise differentiable.

\section{Equivalence of strong and weak forms}

The two formulations are equivalent for smooth data: any strong solution
satisfies the weak form (by construction); conversely, if a weak solution
is smooth enough, reversing the integration by parts and invoking the
arbitrariness of $\vec w$ recovers both the equilibrium PDE (in
$\domain$) and the natural boundary condition (on $\GammaN$). The weak
form is thus not an approximation --- it is a mathematically equivalent
restatement that happens to be far friendlier to discretisation.

\section{The minimum-potential-energy principle}

For elastostatics the weak form has a third face. Define the
\textbf{total potential energy} functional
\begin{equation}
  \Pi(\disp) = \underbrace{\tfrac12\int_\domain \strain(\disp):\Dmat:\strain(\disp)\dd V}_{\text{strain energy }U}
   - \int_\domain \disp\cdot\bodyforce\dd V - \int_{\GammaN}\disp\cdot\bar{\traction}\dd A .
  \label{eq:potential}
\end{equation}
Its first variation in an admissible direction $\vec w\in\mathcal V$ is
\begin{equation}
  \delta\Pi = \int_\domain \strain(\vec w):\Dmat:\strain(\disp)\dd V
   - \int_\domain \vec w\cdot\bodyforce\dd V - \int_{\GammaN}\vec w\cdot\bar{\traction}\dd A,
\end{equation}
which is \emph{exactly} the weak form~\eqref{eq:weak}. Setting
$\delta\Pi=0$ therefore yields the equilibrium solution. Because
$\Dmat$ is symmetric positive-definite, $\Pi$ is strictly convex and the
stationary point is a unique global \textbf{minimum}. Hence:
\begin{equation}
  \text{equilibrium} \iff \delta\Pi=0 \iff \disp \text{ minimises } \Pi .
\end{equation}

\begin{remark}[When the energy view fails]
The equivalence ``Galerkin $\equiv$ energy minimisation'' holds only for
\emph{self-adjoint} problems with a symmetric bilinear form --- linear
elastostatics and steady heat conduction qualify. Non-self-adjoint
problems (e.g.\ convection--diffusion) have no such energy functional; one
must then use the weighted-residual (Galerkin) route directly, which is
why we lead with the weak form rather than the energy principle.
\end{remark}

\section{Weighted residuals and Galerkin's method}

The most general route defines the \textbf{residual}
$\vec R(\disp)=\divergence\stress(\disp)+\bodyforce$ and demands it be
orthogonal to a space of weighting functions:
$\int_\domain \vec w\transpose\vec R\dd V=0$ for all $\vec w$. Choosing
the weighting functions from the \emph{same} space as the trial
functions is the \textbf{(Bubnov--)Galerkin method}; choosing a
different space gives \emph{Petrov--Galerkin}. Integrating by parts as
above turns the weighted-residual statement into the weak
form~\eqref{eq:weak}. For elastostatics all three viewpoints --- virtual
work, energy minimisation, Galerkin --- deliver identical discrete
equations, which we build in \cref{ch:discretization}.
